% =====================================================================
\section{Numerical Experiments}
\label{sec:exp}

\subsection{Scope of the Results}
\label{sec:exp-scope}
The results come from a 30-day continuous diagnostic run with random seed
9{,}900{,}950, which was not used in training. We compare the middle 28
days and use the first and last days only to connect the continuous
boundaries. During evaluation, weights are frozen and exploration is set
to zero. All policies share the same truck, vessel, and service random
streams on a given day. We record the four cost terms, demand level,
transaction types, and unfinished jobs for each day. The source code,
raw policy results, and figure scripts are released with the paper.

\subsection{Confirmation of Architectural Operation}
\label{sec:exp-confirm}
Implementation checks, commitment constraints, and automated tests
confirm four structural properties. Proposal and acceptance use separate
models. One-sided consent never produces a commitment. Realised future
arrival and completion times are excluded from policy inputs. State also
continues across day boundaries. The run logs confirm two additional
properties. Frozen-policy evaluation made \textbf{zero} counterfactual
simulation calls. The identity check also passed at all four sampled
points with counterfactual labels: a branch that changed no action
reproduced the action and cost of the main run.

\subsection{Paired Daily Comparison Across Policies}
\label{sec:exp-paired}
We run every policy with the same random streams to reduce simulation
noise in pairwise differences, following Kleijnen's common-random-number
method~\cite{ref22}. Cost difference is defined as proposed minus
comparison; a negative value means that the proposed policy costs less.
We treat days as independent comparison units and use a two-sided sign
test (Table~\ref{tab:policies}). A Wilcoxon signed-rank test, which also
uses the size of each daily difference, gives the same conclusions. For
every rule-based comparison, $p\le0.017$.

\begin{table}[tb]
\centering
\caption{Policies evaluated on the same 28 days, sorted by cost
reduction. The two rightmost columns compare each policy with the
proposed policy by day. Policies that adjust arrival time reduce cost,
whereas block-only policies remain near zero.}
\label{tab:policies}
\small
\begin{tabular}{@{}llrrr@{}}
\toprule
& & & \multicolumn{2}{c}{vs.\ proposed} \\
\cmidrule(l){4-5}
Policy & Action range & Reduction & Days lower & $p$ \\
\midrule
\textbf{Proposed policy} & block\,+\,time & \textbf{+14.72\%} & --- & --- \\
Proposed, time only & time & +13.44\% & --- & --- \\
Untrained model & block\,+\,time & +7.97\% & 18/28 & 0.185 \\
\midrule
Rule: least-loaded slot & time & +3.41\% & --- & --- \\
Rule: least-loaded block\,+\,slot & block\,+\,time & +1.51\% & 20/28 & \textbf{0.036} \\
\midrule
Rule: least slack & block & +0.53\% & 20/28 & \textbf{0.036} \\
Rule: first-come-first-served & block & +0.17\% & 22/28 & \textbf{0.004} \\
Rule: shortest processing time & block & +0.05\% & 21/28 & \textbf{0.013} \\
Rule: net gain & block & $-$0.48\% & 21/28 & \textbf{0.013} \\
Proposed, block only & block & $-$1.31\% & --- & --- \\
\midrule
No reallocation & --- & 0.00\% & 26/28 & \textbf{$<$0.001} \\
\bottomrule
\end{tabular}
\end{table}

Over 28 days, total cost was KRW 10.00\,bn without reallocation and KRW
8.52\,bn with the proposed policy. The reduction was KRW 1.471\,bn, or
14.7\%. Table~\ref{tab:policies} reports the reduction for each policy.
The four existing rules use block reallocation only, so this comparison
does not separate policy quality from action range. Sect.~\ref{sec:exp-rules}
provides a matched comparison. The proposed policy was more expensive on
only two days, and both differences were small, whereas the largest
reductions all occurred on congested days (Fig.~\ref{fig:paired}).

\begin{center}
\begin{minipage}{\textwidth}
\centering
\includegraphics[width=\textwidth]{final_figure/fig-paired-daily.pdf}
\captionof{figure}{Paired daily cost difference against no reallocation over
the same 28 days with the same random streams. Days are sorted by the
difference, and a positive value means the proposed policy costs less.
(a) All 28 days; four congested days dominate the axis, so (b) enlarges
the shaded ranks 1--24. The proposed policy costs less on 26 of the 28
days (two-sided sign test, $p<0.001$). The two exceptions are small,
KRW $-17.1$\,m and KRW $-0.5$\,m, whereas each of the four congested days
exceeds KRW 290\,m.}
\label{fig:paired}
\end{minipage}
\end{center}

\subsection{Demand Level and the Decomposition of Actions}
\label{sec:exp-decomp}
Table~\ref{tab:actions} decomposes the reduction by demand level. The
four days with 12{,}500 or 15{,}000 trucks account for 90.5\% of the
total. This concentration supports developing the architecture toward a
policy that \emph{intervenes selectively when congestion is observed}
rather than continuously. Restricting the policy to one action type at a
time separates the spatial and temporal contributions. Arrival-time
adjustment accounts for 91.3\% of the total reduction, mostly on the four
congested days (Fig.~\ref{fig:decomp}). Block-only reallocation saves KRW 0.124\,bn under smooth
and normal demand but raises cost at the highest demand, producing a net
increase of KRW 0.131\,bn. The policy should therefore adjust its use of
spatial and temporal actions by demand level. Time-slot quotas were open
in this diagnostic run. The arrival-time results thus describe the
\emph{potential range} of time shifting, not an executable schedule.
Sensitivity analysis with finite quotas can define the practical range.

\begin{table}[tb]
\centering
\caption{Cost reduction by demand level (KRW bn, relative to no
reallocation). Columns three to five separate the allowed actions. The
last column compares the trained and untrained models. Arrival-time
adjustment and learning contribute mainly on the four congested days,
whereas block reallocation raises cost at the highest demand.}
\label{tab:actions}
\small
\begin{tabular}{@{}rrrrrr@{}}
\toprule
Daily & & Both & Time & Block & Due to \\
demand & Days & actions & only & only & training \\
\midrule
3,500 & 12 & 0.047 & $-$0.013 & 0.057 & 0.056 \\
5,000 & 8 & 0.060 & 0.029 & 0.067 & 0.029 \\
7,500 & 4 & 0.032 & 0.008 & 0.000 & $-$0.049 \\
12,500 & 2 & 0.611 & 0.573 & $-$0.014 & 0.144 \\
15,000 & 2 & 0.722 & 0.747 & $-$0.241 & 0.495 \\
\midrule
\textbf{Total} & \textbf{28} & \textbf{1.471} & \textbf{1.344} & \textbf{$-$0.131} & \textbf{0.675} \\
\multicolumn{2}{@{}l}{\textit{share of total}} & 100\% & \textbf{91.3\%} & --- & \textbf{45.9\%} \\
\bottomrule
\end{tabular}
\end{table}

\begin{center}
\begin{minipage}{\textwidth}
\centering
\includegraphics[width=\textwidth]{final_figure/fig-decomposition.pdf}
\captionof{figure}{Cost reduction by demand level, decomposed by action
range; the values are those of Table~\ref{tab:actions}. The vertical axis
is split by demand because the two congested levels are more than ten
times larger than the rest: (a) 3{,}500--7{,}500 and (b)
12{,}500--15{,}000 trucks per day, with the number of days observed
printed under each tick. The first three bars are measured against no
reallocation. The fourth bar, set slightly apart, uses a different
baseline: it is the trained model minus the untrained one, that is, the
last column of Table~\ref{tab:actions}. Arrival-time adjustment carries
91.3\% of the total reduction, whereas block-only reallocation turns
negative once demand is congested.}
\label{fig:decomp}
\end{minipage}
\end{center}

\subsection{Rule-Based Policies of the Same Action Range}
\label{sec:exp-rules}
To determine whether the difference in Sect.~\ref{sec:exp-decomp} comes
from the action itself, the rule-based policies must receive the same
action range. We therefore transfer the \emph{least-loaded} criterion
from the spatial axis to the temporal axis and add two rules. This change
introduces no new parameter and uses the same congestion trigger as the
other rules. Three findings follow. First, arrival-time adjustment under
a rule reduces cost by 3.41\% relative to no reallocation
($p=0.0125$). The temporal action therefore has value by itself. Second,
with the full spatio-temporal action range matched, the proposed policy
costs less than the rule on 20 of 28 days ($p=0.036$). Third, for
arrival-time adjustment alone, the rule costs less on more days; the
proposed policy is lower on only 8 of 28 days. Yet the proposed policy's
total reduction is about four times larger. The rule is slightly better
under smooth demand, whereas the proposed policy reduces cost 3.8--5.6
times more on congested days. Thus, the learned policy tends to produce
large savings on expensive days rather than frequent small wins. Both
the number of wins and the total reduction are needed to show this
pattern. The rule using both actions also weakens under congestion. This
suggests that the congestion sensitivity of block reallocation may be a
property of the action, not only of one policy.

\subsection{Learning Effect and Dispatching Sensitivity}
\label{sec:exp-learning}
We first examine model fit. Across 30 training iterations, the proposal
network's Huber loss decreases from 0.079 to 0.029 on training decisions
and from 0.268 to 0.119 on validation decisions. The acceptance
network's loss decreases from 0.171 to 0.027 and from 0.315 to 0.136,
respectively. These values are means over the first and last five
iterations, with targets in units of KRW 100{,}000. The decline in
validation loss indicates that the fit extends to unseen decisions.
However, final validation loss remains about four times training loss.
With about 56 labelled decisions per iteration, this gap supports the
choice of a small model. As exploration decreases, the share of
decisions with exactly zero counterfactual difference falls from 51.4\%
to 37.6\%. Training runs 4{,}671 counterfactual worlds in total.
Fig.~\ref{fig:learning} shows both loss curves together with the
exploration schedule that accompanies them.

\begin{center}
\begin{minipage}{\textwidth}
\centering
\includegraphics[width=\textwidth]{final_figure/fig-learning-curve.pdf}
\captionof{figure}{Training and validation loss over 30 iterations, with the
per-iteration value drawn faintly beneath a five-iteration mean. (a) The
proposal network and (b) the acceptance network are drawn separately
because they are fitted on different decisions, and the vertical axis is
logarithmic because the loss spans three orders of magnitude. (c) The
exploration probability $\varepsilon$ decays linearly from 0.50 to 0.05.
Shading marks the iterations whose demand level was drawn at 12{,}500
trucks per day or above. Because the demand level is redrawn every
iteration, much of the visible fluctuation, including the rise around
iterations 17--24, follows that draw rather than the fit.}
\label{fig:learning}
\end{minipage}
\end{center}

Cost is KRW 10.00\,bn without reallocation, KRW 9.20\,bn with the
untrained model, and KRW 8.52\,bn with the trained model. Training
therefore adds KRW 0.675\,bn, or 45.9\% of the total reduction. However,
the trained model costs less on only 18 of 28 days ($p=0.1849$). The
mean daily difference is KRW $-24.1$\,m, about 18 times the median of
$-1.3$\,m. This gap shows that the learning effect is concentrated rather
than evenly distributed. The four congested days account for 94.6\% of
the learning effect, and the trained model costs less on all four. The
corresponding counts at lower demand levels are 9/12, 5/8, and 0/4. For
the remaining 24 days, the mean difference is $-1.5$\,m and the
bootstrap interval includes zero.

Resampling daily differences within demand strata gives a 95\% interval
of $[-32,-17]$\,m per day for the mean difference. This procedure
preserves the demand mix fixed by the design. Without stratification, the
interval widens to $[-52,-3]$\,m, and 99.2\% of resamples remain
negative. Both intervals exclude zero, but the stratified interval is
partly narrow by construction because each congested stratum contains
only two days. The interpretation is therefore limited. Learning
produces a large total reduction and is lower on all four congested
days, but the day-level sign test is not significant. Establishing the
learning effect at a conventional significance level requires more
congested days, not simply more days overall.

The ranking of crane-dispatching rules also changes with demand. The
shortest processing time rule is 22--25\% more expensive than the other
rules under smooth and normal demand, but it is the least expensive from
7{,}500 trucks upward. It also has the lowest overall cost under the
designed demand mix. The effect of reallocation should therefore be
tested under other dispatching rules.
